\section{Konklusion}
Det kan konkluderes, at gruppen har haft en overordnet effektiv og velfungerende projektproces, hvor en stor del af de planlagte story points er blevet gennemført på trods af tekniske, organisatoriske og samarbejdsmæssige udfordringer undervejs i forløbet. Projektet har udviklet sig fra en tidlig idéfase til et fungerende multiplayer-spil med implementerede kerne-systemer såsom sessionshåndtering, team- og klassesystem, AI, win condition og et sammenhængende gameplay-loop. Gruppen har undervejs opnået erfaring med både teknisk spiludvikling i Unreal Engine samt struktureret projektarbejde gennem en Scrum-lignende arbejdsform.

En væsentlig læring fra projektforløbet er, at en klar rollefordeling og en struktureret arbejdsproces har haft stor betydning for fremdriften. Det har været en styrke, at gruppen tidligt etablerede roller som Scrum Master, Product Owner og analytisk ansvarlig, da dette har skabt overblik og reduceret risikoen for ansvarsforvirring. Samtidig har brugen af sprintbaseret udvikling og løbende evalueringer gjort det muligt at tilpasse projektets retning undervejs.

Samtidig har projektet også vist, at der er områder, hvor gruppen fremadrettet kan forbedre deres arbejdsmetode. Manglende daglige stand-ups har i nogle tilfælde betydet, at mindre problemer ikke blev identificeret tidligt nok. Derudover har visse tekniske beslutninger, såsom tidlig implementering af komplekse systemer uden fuld afklaring, medført ekstra arbejde senere i forløbet.

På baggrund af disse erfaringer kan der opstilles følgende anbefalinger til fremtidigt gruppearbejde:

Anbefalinger (do’s and don’ts):

DO: Etabler klare roller tidligt i projektet og fasthold dem gennem hele forløbet.
DO: Arbejd i korte og veldefinerede iterationer med tydelige mål for hvert sprint.
DO: Brug løbende evaluering (retrospektiver) aktivt til at forbedre arbejdsprocessen.
DO: Dokumentér beslutninger løbende for at undgå misforståelser senere i projektet.
DO: Vær proaktiv i kommunikationen, både internt i gruppen og med vejleder.
DON’T: Undlad at vente med at adressere tekniske eller organisatoriske problemer.
DON’T: Undervurdér kompleksiteten i tidlige systemer (især multiplayer og AI).
DON’T: Undervurdér kompleksiteten i VoiceChat Systemer.
DON’T: Arbejd isoleret uden løbende koordinering, da det kan skabe integrationsproblemer.
DON’T: Spring struktureret test og review over, da det ofte fører til fejl senere i processen.

Samlet set har projektet været en lærerig proces, hvor gruppen har opnået både tekniske og samarbejdsmæssige kompetencer. Kombinationen af en Scrum-lignende metode, tydelig rollefordeling og løbende iterationer har været afgørende for, at projektet kunne gennemføres med et funktionelt slutprodukt. Samtidig har udfordringerne undervejs bidraget til værdifulde erfaringer, som gruppen kan anvende i fremtidige projekter for at forbedre både effektivitet og kvalitet i samarbejdet.